\documentclass[12pt,a4paper]{article}
\usepackage[UTF8]{ctex}
\usepackage{amsmath,amssymb,amsthm}
\usepackage{geometry}
\usepackage{bm}
\usepackage{hyperref}
\usepackage{tikz}
\usetikzlibrary{arrows,arrows.meta,positioning,calc,shapes.geometric}

\geometry{margin=2.5cm}

\title{例题12.5的详细理解：多个接触的CoR表示}
\author{第12章抓取与操作补充说明}
\date{}

\newtheorem{theorem}{定理}
\newtheorem{definition}{定义}
\newtheorem{example}{例子}

\begin{document}

\maketitle

\section{问题1：灰色图是CoR点的合集吗？}

\textbf{答案}：是的！灰色区域表示\textbf{可行CoR的集合}。

\begin{itemize}
\item 灰色区域内的每个点都对应一个可行的CoR
\item 每个CoR对应一个可行的twist（运动）
\item 灰色区域是\textbf{凸的}（convex）
\item 这是多个接触约束的交集
\end{itemize}

\section{例题12.5的设置}

\begin{example}[例题12.5]
图12.12(a)显示一个平面物体站在桌子上，同时被一个静止的机器人手指接触。手指定义了一个不等式约束，桌子定义了两个约束。不违反不可穿透约束的twist锥由对每个接触一致标记的CoR表示（图12.12(b)）。每个可行CoR都标记有一个接触模式，该模式是各个接触标签的连接（图12.12(c)）。
\end{example}

\subsection{接触设置}

\begin{itemize}
\item \textbf{接触1}：桌子提供的接触约束1
\item \textbf{接触2}：桌子提供的接触约束2
\item \textbf{接触3}：静止机器人手指提供的接触约束
\item 总共：\textbf{3个接触约束}
\end{itemize}

\subsection{约束类型}

所有接触都是\textbf{静止的}（stationary），因此：
\begin{itemize}
\item 约束是\textbf{齐次的}
\item 可行twist集合是\textbf{以原点为根部的凸锥}
\item 可以用CoR表示
\end{itemize}

\section{图12.12(b)：可行CoR区域（灰色区域）}

\subsection{灰色区域的含义}

\textbf{灰色区域}表示所有\textbf{可行CoR的集合}：

\begin{itemize}
\item 区域内的每个点都对应一个可行的CoR
\item 每个CoR对应一个可行的twist（运动）
\item 这些twist满足所有3个接触的不可穿透约束
\end{itemize}

\subsection{为什么是凸的？}

\begin{itemize}
\item 每个接触贡献一个约束区域（半空间）
\item 可行CoR区域是所有约束区域的\textbf{交集}
\item 半空间的交集是\textbf{凸集}
\item 因此灰色区域是凸的
\end{itemize}

\subsection{边界和无穷远}

\textbf{重要观察}：根据图12.12的说明，延伸到左侧和底部的线在无穷远处"环绕"，从右侧和顶部回来。

这意味着：
\begin{itemize}
\item 在射影几何中，这是一个\textbf{单连通的凸区域}
\item 包括无穷远点（对应平移）
\item 在有限平面上，我们看到的是延伸到边界的区域
\end{itemize}

\section{图12.12(c)：接触模式标记}

\subsection{接触模式的含义}

每个可行CoR都标记有一个\textbf{接触模式}，该模式是各个接触标签的连接。

\textbf{例子}：
\begin{itemize}
\item \textbf{RRR}：所有三个接触都是滚动的（物体完全静止）
\item \textbf{BSrB}：接触1分离，接触2向右滑动，接触3分离
\item \textbf{SrBSr}：接触1向右滑动，接触2分离，接触3向右滑动
\end{itemize}

\subsection{接触模式的分布}

\begin{itemize}
\item \textbf{内部区域}：所有接触都分离（BBB）
\item \textbf{边界线}：一个接触约束主动（如BSrB、BSlB等）
\item \textbf{顶点/交点}：两个或更多接触约束主动（如RRR、SrSrSr等）
\end{itemize}

\section{关键观察：一阶分析的局限性}

\subsection{问题：$S_rBS_r$接触模式}

图12.12(c)中标记了$S_rBS_r$接触模式，但原书指出：

\textbf{这个运动实际上是不可能的！}

\textbf{原因}：
\begin{itemize}
\item 物体会立即穿透静止的手指
\item 一阶分析忽略了局部接触曲率
\item 二阶分析会显示这个运动是不可能的
\end{itemize}

\subsection{一阶 vs 二阶分析}

\textbf{一阶分析}：
\begin{itemize}
\item 只考虑接触法线方向
\item 忽略接触曲率
\item 可能错误地认为某些运动是可行的
\end{itemize}

\textbf{二阶分析}：
\begin{itemize}
\item 考虑接触曲率$\partial^2 d/\partial q^2$
\item 更准确地判断运动的可行性
\item 可能将一阶的滚动-滑动运动重新分类为分离或穿透
\end{itemize}

\subsection{特殊情况}

\textbf{注意}：如果物体在接触处的曲率半径足够小，那么$S_rBS_r$运动可能是可行的。

\section{详细分析：三个接触的约束}

\subsection{接触1：桌子约束1}

\begin{itemize}
\item 接触法线：$\hat{n}_1$（指向物体内部）
\item 约束：$F_1^T V \geq 0$
\item CoR标记规则：
  \begin{itemize}
  \item 法线左侧：'+'
  \item 法线右侧：'-'
  \item 法线上：'$\pm$'
  \end{itemize}
\end{itemize}

\subsection{接触2：桌子约束2}

\begin{itemize}
\item 接触法线：$\hat{n}_2$（指向物体内部）
\item 约束：$F_2^T V \geq 0$
\item CoR标记规则：同上
\end{itemize}

\subsection{接触3：手指约束}

\begin{itemize}
\item 接触法线：$\hat{n}_3$（指向物体内部）
\item 约束：$F_3^T V \geq 0$
\item CoR标记规则：同上
\end{itemize}

\subsection{可行CoR区域}

可行CoR区域是三个约束区域的\textbf{交集}：

\begin{equation}
\text{可行CoR} = \{\text{CoR} | \text{CoR满足所有三个接触的约束}\}
\end{equation}

\textbf{关键}：
\begin{itemize}
\item 每个接触贡献一个约束区域
\item 可行区域是这些区域的交集
\item 交集是凸的
\end{itemize}

\section{接触模式的确定}

\subsection{对于每个可行CoR}

对于可行CoR区域内的每个点，我们需要确定：

\begin{enumerate}
\item \textbf{每个接触的状态}：
   \begin{itemize}
   \item CoR在接触法线左侧：'+'标记 → 接触分离（B）
   \item CoR在接触法线右侧：'-'标记 → 接触分离（B）
   \item CoR在接触法线上：'$\pm$'标记 → 接触可能是滚动（R）或滑动（S）
   \end{itemize}

\item \textbf{如果是滑动}：
   \begin{itemize}
   \item 需要进一步判断是$S_l$还是$S_r$
   \item 这取决于旋转方向和CoR位置
   \end{itemize}

\item \textbf{接触模式}：
   \begin{itemize}
   \item 连接三个接触的标签
   \item 例如：$S_rBS_r$、$RRR$、$BBB$等
   \end{itemize}
\end{enumerate}

\section{具体例子分析}

\subsection{例子1：CoR在内部（BBB）}

\begin{itemize}
\item CoR不在任何接触法线上
\item 所有接触都分离
\item 接触模式：\textbf{BBB}
\end{itemize}

\subsection{例子2：CoR在一条边界上（BSrB）}

\begin{itemize}
\item CoR在接触2的法线上
\item 接触1和3分离（B）
\item 接触2滑动（$S_r$）
\item 接触模式：\textbf{BSrB}
\end{itemize}

\subsection{例子3：CoR在顶点（RRR）}

\begin{itemize}
\item CoR在所有三条接触法线的交点上
\item 所有接触都滚动
\item 接触模式：\textbf{RRR}
\item 这对应零速度（物体完全静止）
\end{itemize}

\section{一阶分析的局限性总结}

\subsection{一阶分析的正确性}

\begin{itemize}
\item \textbf{总是正确}：分离（B）和穿透（不可行）
\item 一阶分析总是能正确识别这些运动
\end{itemize}

\subsection{一阶分析的不确定性}

\begin{itemize}
\item \textbf{可能错误}：滚动-滑动运动（R或S）
\item 一阶分析认为可行的运动，二阶分析可能显示为分离或穿透
\item 例如：$S_rBS_r$在一阶分析中可行，但二阶分析显示不可行
\end{itemize}

\subsection{形状闭合的保证}

\begin{itemize}
\item 如果一阶分析显示形状闭合，则\textbf{任何分析}都显示形状闭合
\item 如果只有滚动-滑动运动在一阶分析中可行，则：
  \begin{itemize}
  \item 可能通过高阶分析实现形状闭合
  \item 或者，在任何分析中都不处于形状闭合
  \end{itemize}
\end{itemize}

\section{总结}

\subsection{关键点}

\begin{enumerate}
\item \textbf{灰色区域}：表示可行CoR的集合（凸集）

\item \textbf{接触模式}：每个CoR对应一个接触模式（三个接触标签的连接）

\item \textbf{一阶分析的局限性}：
   \begin{itemize}
   \item 可能错误地认为某些运动可行
   \item 需要二阶或更高阶分析来验证
   \end{itemize}

\item \textbf{形状闭合}：
   \begin{itemize}
   \item 如果一阶分析显示形状闭合，则保证形状闭合
   \item 否则需要更高阶分析
   \end{itemize}
\end{enumerate}

\subsection{应用}

这个例题展示了：
\begin{itemize}
\item 如何用CoR表示多个接触的约束
\item 如何确定接触模式
\item 一阶分析的优点和局限性
\end{itemize}

\end{document}

